% GENERATED FILE. Edit canonical lessons, then run npm run generate:books.
% canonical-source-hash: fnv1a64:0f9c7a29f8889afe
% canonical-lessons: FR-C08-heure, FR-C08-il-est, FR-C08-une-heure, FR-C08-midi, FR-C08-minuit, FR-C08-quart, FR-C08-demie, FR-C08-moins, FR-C08-practice

\chapter{Telling the Time}
\label{ch:time}
\hlchaptermodality{8}

\begin{chapteropening}
\textbf{By the end of this chapter:} \emph{I can ask for and tell the time in French.}

One noun, one frame and three small additions. heure is the same word as English hour, both from Latin hora and both keeping an h neither language ever said; il est ... heures counts the hours where English says o'clock; and midi and minuit are one idea used twice, the middle of the day and the middle of the night.
\end{chapteropening}

\section[heure]{heure — the hour}

\label{lesson:FR-C08-heure}



\begin{quote}
With your numbers in hand you are one word away from telling the time. This is that word — and it turns out you have been saying it in English your whole life.
\end{quote}

\begin{sounds}
\begin{itemize}
\raggedright
  \item \texttt{silent-h} — the \textbf{h} of \textbf{heure} is silent, exactly as in English \emph{hour}. Say \emph{eur}.
  \item \texttt{liaison-z} — because the word \textbf{begins with a vowel sound}, the number in front links straight onto it. That is the liaison you met counting: \textbf{deux heures} = \emph{deu-\textbf{z}-eur}, \textbf{trois heures} = \emph{troi-\textbf{z}-eur}.
\end{itemize}

The silent \emph{h} is why the liaison happens at all. A written \emph{h} that is never pronounced is not a consonant in the way; the number hears a vowel and reaches for it.
\end{sounds}

\begin{cousinweb}
\textbf{heure} comes from Latin \textbf{hōra}, which came from Greek \textbf{hōrā}. The Greek word did not mean \textquotedblleft{}hour.\textquotedblright{} It meant \emph{a period, a season, the right time} — the goddesses of the seasons were the \textbf{Horae}. Latin narrowed it to a fixed slice of the day, and French wore \emph{hōra} down to \textbf{heure}.

English took the \textbf{same} Latin \emph{hōra} and made \textbf{hour}. So \emph{heure} and \emph{hour} are one word spelt two ways — and the silent \emph{h} in both is a fossil of the Latin spelling that neither language ever said aloud.
\end{cousinweb}

\subsection*{Your turn}
Say these aloud:

\begin{itemize}
\raggedright
  \item \textquotedblleft{}heure\textquotedblright{} — \emph{eur}, no \emph{h} at all
  \item \textquotedblleft{}deux heures\textquotedblright{} — \emph{deu-z-eur}, the number linking on
  \item \textquotedblleft{}trois heures, six heures, dix heures\textquotedblright{} — hear the \emph{z} each time
\end{itemize}

\begin{tcolorbox}[breakable,colback=teal!4,colframe=teal!35!black,title={Before you move on}]
How do you say hour? (\textbf{heure}, \emph{eur}.) Why does a number link onto it with a \emph{z}? (The \textbf{h} is silent, so the word starts with a vowel sound.) What did the Greek \emph{hōrā} originally mean? (\textbf{A season, the right time} — the \emph{Horae} were the season-goddesses.) Next: the frame that puts a number in front of it.
\end{tcolorbox}
\section[il est … heures]{il est … heures — telling the clock}

\label{lesson:FR-C08-il-est}



\begin{quote}
One word for the hour, one small frame to hang it in, and the clock is yours. Take the frame whole for now — it never changes.
\end{quote}

\begin{grammarlens}[title={Grammar Lens: French counts the hours}]
To say the time, French uses a fixed opening — \textbf{il est} — then the number, then \textbf{heures}:

\begin{quote}
\textbf{il est} + \emph{number} + \textbf{heures}
\end{quote}

\textbf{Il est deux heures.} Read it literally and it says \emph{\textquotedblleft{}it is two \textbf{hours}.\textquotedblright{}} French \textbf{counts the hours}. English does something stranger and then hides it: \emph{two o'clock} is a worn-down \emph{two of the clock}, which counted the hours too before it shrank.

Learn \textbf{il est} as one unbreakable piece for now. It is the same \emph{it is} that English uses for weather and time — a subject that stands for nothing in particular — and you will meet the verb inside it properly later.
\end{grammarlens}

\begin{sounds}
\begin{itemize}
\raggedright
  \item \textbf{il est} = \emph{eel-\textbf{eh}}. The final \textbf{t} of \emph{est} is silent on its own.
  \item But before \textbf{heures}, with its silent \emph{h}, that \textbf{t} wakes up and links: \textbf{il est une heure} = \emph{eel eh-\textbf{t}-ün-eur}. Another liaison, the same instinct as \emph{deux heures}.
\end{itemize}
\end{sounds}

\subsection*{Your turn}
Say these aloud:

\begin{itemize}
\raggedright
  \item \textquotedblleft{}il est\textquotedblright{} — \emph{eel-eh}, no \emph{t} sounded
  \item \textquotedblleft{}Il est deux heures.\textquotedblright{} — \emph{eel eh deu-z-eur}
  \item \textquotedblleft{}Il est six heures. Il est dix heures.\textquotedblright{}
\end{itemize}

\emph{Twice through:} \textquotedblleft{}Il est deux heures.\textquotedblright{}

\begin{tcolorbox}[breakable,colback=teal!4,colframe=teal!35!black,title={Before you move on}]
What frame tells the time in French? (\textbf{il est} + number + \textbf{heures}.) What does \emph{Il est deux heures} say literally? (\textquotedblleft{}It is two \textbf{hours}.\textquotedblright{}) What is English \emph{o'clock} worn down from? (\emph{Of the clock} — which counted hours too.) Next: the one hour that takes a different ending.
\end{tcolorbox}
\section[une heure]{une heure — one o'clock}

\label{lesson:FR-C08-une-heure}



\begin{quote}
Eleven of the twelve hours behave identically. One does not, and the reason is a rule you already met on a different word.
\end{quote}

\begin{grammarlens}[title={Grammar Lens: one hour, two hours}]
\textbf{heure} is a \textbf{feminine} noun, so \textquotedblleft{}one\textquotedblright{} in front of it is \textbf{une}, not \emph{un} — the same gendered \textquotedblleft{}one\textquotedblright{} you learned counting. And because the noun counts, it takes a plural \textbf{-s} from two upwards:

\begin{itemize}
\raggedright
  \item \textbf{Il est une heure.} — one o'clock. One hour, so no \emph{-s}.
  \item \textbf{Il est deux heures.} — two o'clock. Two hours, so \emph{-s}.
\end{itemize}

Now listen rather than read. That \textbf{-s} is \textbf{silent}, like every plural \emph{-s} in French. \emph{Une heure} and \emph{deux heures} end in the identical sound. The plural is visible on the page and inaudible in the air — what your ear actually uses to tell them apart is the \textbf{number in front}, which is why French never drops it.
\end{grammarlens}

\subsection*{Your turn}
Say these aloud:

\begin{itemize}
\raggedright
  \item \textquotedblleft{}Il est une heure.\textquotedblright{} — \emph{eel eh-t-ün-eur}
  \item \textquotedblleft{}Il est deux heures.\textquotedblright{} — \emph{eel eh deu-z-eur}
  \item both again and notice the ending never changes — only the number does
\end{itemize}

\begin{tcolorbox}[breakable,colback=teal!4,colframe=teal!35!black,title={Before you move on}]
How do you say one o'clock? (\textbf{Il est une heure} — \emph{une}, because \emph{heure} is feminine.) When does \emph{heure} take an \textbf{-s}? (From \textbf{two} upwards.) Can you hear that \emph{-s}? (\textbf{No} — it is silent, so the number carries the count.) Next: the two hours that refuse a number altogether.
\end{tcolorbox}
\section[midi]{midi — noon}

\label{lesson:FR-C08-midi}



\begin{quote}
Two hours in the day do not get a number. They get a name, because they are the pivots the rest of the day turns on. Here is the first.
\end{quote}

\begin{sounds}
\begin{itemize}
\raggedright
  \item \textbf{midi} = \emph{mee-DEE}, two clean syllables with the weight at the end.
  \item Both vowels are the tight French \textbf{i} of \emph{machine}, never the \emph{i} of \emph{sit}.
\end{itemize}
\end{sounds}

\begin{cousinweb}
\textbf{midi} is Latin \textbf{medius diēs}, \textquotedblleft{}the middle day,\textquotedblright{} and both halves are still visible:

\begin{quote}
\textbf{mi-} (\textquotedblleft{}middle,\textquotedblright{} from \emph{medius}) + \textbf{-di} (\emph{diēs}, \textquotedblleft{}day\textquotedblright{}) = \textbf{the day's middle}
\end{quote}

You already own the second half. That \textbf{-di} is the same \emph{diēs} that ends \emph{lundi} and \emph{mardi} — here it has come round to the front of a different word. And \emph{medius} is the English \textbf{medium}, \textbf{median}, \textbf{mid} and \textbf{middle}, all one family.
\end{cousinweb}

\begin{grammarlens}[title={Grammar Lens: a name instead of a number}]
Because \textbf{midi} \emph{is} the hour, it does not take \textbf{heures}. Drop the number and drop the noun:

\begin{quote}
\textbf{Il est midi.} — \textquotedblleft{}It is noon.\textquotedblright{}
\end{quote}

Not \emph{il est douze heures}. The frame stays; what goes into it is a name.
\end{grammarlens}

\begin{culture}
English \textbf{noon} is Latin \textbf{nōna hōra}, \textquotedblleft{}the \textbf{ninth} hour\textquotedblright{} — which in Roman reckoning was the middle of the afternoon. Over centuries the meal it named crept earlier, and the word came with it until it landed on midday and stopped. So English says \emph{noon} and means \textquotedblleft{}the ninth hour\textquotedblright{}; French says \emph{midi} and means what it says.
\end{culture}

\subsection*{Your turn}
Say these aloud:

\begin{itemize}
\raggedright
  \item \textquotedblleft{}midi\textquotedblright{} — \emph{mee-DEE}
  \item \textquotedblleft{}Il est midi.\textquotedblright{} — no number, no \emph{heures}
  \item it split — \textquotedblleft{}mi- plus -di: the day's middle\textquotedblright{}
\end{itemize}

\begin{tcolorbox}[breakable,colback=teal!4,colframe=teal!35!black,title={Before you move on}]
How do you say noon? (\textbf{midi}, \emph{mee-DEE}.) What are its two halves? (\textbf{mi-}, \textquotedblleft{}middle,\textquotedblright{} and \textbf{-di}, \emph{diēs}, \textquotedblleft{}day.\textquotedblright{}) Where have you met that \emph{-di} before? (Ending \textbf{lundi} and \textbf{mardi}.) Why is there no \emph{heures} after it? (\emph{Midi} \textbf{is} the hour.) Next: its opposite number.
\end{tcolorbox}
\section[minuit]{minuit — midnight}

\label{lesson:FR-C08-minuit}



\begin{quote}
The second pivot, and it is built from a word you have had since the very first chapter.
\end{quote}

\begin{sounds}
\begin{itemize}
\raggedright
  \item \textbf{minuit} = \emph{mee-NWEE}. The \textbf{ui} is a single gliding sound, \emph{wee}, not two vowels taken in turn.
  \item The final \textbf{t} is silent, as it is in \textbf{nuit} itself.
\end{itemize}
\end{sounds}

\begin{cousinweb}
\textbf{minuit} is Latin \textbf{media nox}, \textquotedblleft{}the middle night,\textquotedblright{} and the two halves are as plain as they were in \emph{midi}:

\begin{quote}
\textbf{mi-} (\textquotedblleft{}middle,\textquotedblright{} \emph{medius} again) + \textbf{-nuit} (\emph{nox}, \textquotedblleft{}night\textquotedblright{}) = \textbf{the night's middle}
\end{quote}

The second half is the \textbf{nuit} you learned in your first week. Latin \emph{nox} has the stem \emph{noct-}, which English carries in \textbf{nocturnal} and, straight through its Germanic side, in \textbf{night} itself.

So the two pivots of the French day are one idea applied twice: the \textbf{middle} of the day, and the \textbf{middle} of the night. English builds \emph{midday} and \emph{midnight} out of exactly the same two pieces.
\end{cousinweb}

\begin{grammarlens}[title={Grammar Lens: the same rule as midi}]
\textbf{minuit} is the hour, so no number and no \textbf{heures}:

\begin{quote}
\textbf{Il est minuit.} — \textquotedblleft{}It is midnight.\textquotedblright{}
\end{quote}
\end{grammarlens}

\subsection*{Your turn}
Say these aloud:

\begin{itemize}
\raggedright
  \item \textquotedblleft{}minuit\textquotedblright{} — \emph{mee-NWEE}, the \emph{ui} as one glide
  \item \textquotedblleft{}Il est midi. Il est minuit.\textquotedblright{} — the two pivots
  \item both split — \textquotedblleft{}the day's middle, the night's middle\textquotedblright{}
\end{itemize}

\begin{tcolorbox}[breakable,colback=teal!4,colframe=teal!35!black,title={Before you move on}]
How do you say midnight? (\textbf{minuit}, \emph{mee-NWEE}.) What are its two halves? (\textbf{mi-}, \textquotedblleft{}middle,\textquotedblright{} and \textbf{nuit}, \textquotedblleft{}night,\textquotedblright{} from Latin \emph{nox}.) Which English word carries the \emph{noct-} stem? (\textbf{Nocturnal}.) Say both pivots. (\emph{Il est midi. Il est minuit.}) Next: the quarter hour.
\end{tcolorbox}
\section[et quart]{et quart — quarter past}

\label{lesson:FR-C08-quart}



\begin{quote}
Whole hours are not enough to make an appointment. Three small additions finish the clock, and the first is already half-known to you from counting.
\end{quote}

\begin{sounds}
\begin{itemize}
\raggedright
  \item \textbf{quart} = \emph{kar}. The \textbf{qu-} is a plain \emph{k}, and the final \textbf{t} is silent.
  \item The \textbf{r} is the throat-\textbf{r}, as in \textbf{quatre}.
\end{itemize}
\end{sounds}

\begin{cousinweb}
\textbf{quart} is \textquotedblleft{}a fourth part,\textquotedblright{} and it comes from the same Latin \textbf{quattuor} that gave you \textbf{quatre}. English took the identical word and made \textbf{quarter} — a fourth of an hour, a fourth of a gallon, a fourth of a town.

Once you see \emph{quatre} inside \emph{quart}, it stops being a new word and becomes a number you already have, doing a different job.
\end{cousinweb}

\begin{grammarlens}[title={Grammar Lens: adding to the hour}]
French adds the quarter with \textbf{et} (\textquotedblleft{}and\textquotedblright{}), after the hour:

\begin{quote}
\textbf{Il est deux heures et quart.} — \textquotedblleft{}It is a quarter past two.\textquotedblright{}
\end{quote}

Read literally: \emph{\textquotedblleft{}it is two hours \textbf{and a quarter}.\textquotedblright{}} French does not say \textquotedblleft{}past\textquotedblright{}; it adds.
\end{grammarlens}

\subsection*{Your turn}
Say these aloud:

\begin{itemize}
\raggedright
  \item \textquotedblleft{}quart\textquotedblright{} — \emph{kar}, no \emph{t}
  \item \textquotedblleft{}Il est deux heures et quart.\textquotedblright{}
  \item \textquotedblleft{}quatre, quart\textquotedblright{} — the number and the fourth part, one root
\end{itemize}

\begin{tcolorbox}[breakable,colback=teal!4,colframe=teal!35!black,title={Before you move on}]
How do you say a quarter? (\textbf{quart}, \emph{kar}.) Which number is it built on? (\textbf{quatre} — Latin \emph{quattuor}.) How do you say quarter past two? (\emph{Il est deux heures \textbf{et quart}.}) Next: the half hour.
\end{tcolorbox}
\section[et demie]{et demie — half past}

\label{lesson:FR-C08-demie}



\begin{quote}
Same move as the quarter, a different fraction — and a spelling that tells you something about the word it leans on.
\end{quote}

\begin{sounds}
\begin{itemize}
\raggedright
  \item \textbf{demie} = \emph{duh-MEE}. The first \textbf{e} is the light, almost-swallowed vowel you met in the middle of \emph{samedi}.
  \item The final \textbf{-e} is silent, so \emph{demie} and \emph{demi} sound identical.
\end{itemize}
\end{sounds}

\begin{cousinweb}
\textbf{demie} comes from Latin \textbf{dimidius}, \textquotedblleft{}halved\textquotedblright{} — literally \emph{di-} (\textquotedblleft{}apart\textquotedblright{}) plus \emph{medius}, the \textquotedblleft{}middle\textquotedblright{} you have now met twice, in \textbf{midi} and \textbf{minuit}. To halve a thing is to cut it at its middle.

English borrowed the French form as the prefix \textbf{demi-}: a \textbf{demigod} is half a god, a \textbf{demitasse} is half a cup.
\end{cousinweb}

\begin{grammarlens}[title={Grammar Lens: why it is demie and not demi}]
The \textbf{-e} on the end is agreement. It follows \textbf{heure}, which is feminine, so the half agrees with it and takes a feminine \textbf{-e}:

\begin{quote}
\textbf{Il est deux heures et demie.} — \textquotedblleft{}It is half past two.\textquotedblright{}
\end{quote}

After \textbf{midi} and \textbf{minuit}, which are masculine, it drops the \emph{-e}:

\begin{quote}
\textbf{Il est midi et demi.} — \textquotedblleft{}It is half past twelve.\textquotedblright{}
\end{quote}

Your ear cannot tell them apart; the page can. That is the same silent-agreement pattern as the plural \emph{-s} on \emph{heures}.
\end{grammarlens}

\subsection*{Your turn}
Say these aloud:

\begin{itemize}
\raggedright
  \item \textquotedblleft{}demie\textquotedblright{} — \emph{duh-MEE}
  \item \textquotedblleft{}Il est deux heures et demie.\textquotedblright{}
  \item \textquotedblleft{}Il est midi et demi.\textquotedblright{} — masculine, so no \emph{-e} on the page
\end{itemize}

\begin{tcolorbox}[breakable,colback=teal!4,colframe=teal!35!black,title={Before you move on}]
How do you say half past two? (\emph{Il est deux heures \textbf{et demie}.}) Why does \emph{demie} carry an \textbf{-e}? (It agrees with \textbf{heure}, which is feminine.) What is inside Latin \emph{dimidius}? (\emph{Medius} — the same \textquotedblleft{}middle\textquotedblright{} as in \emph{midi}.) Next: the one that subtracts.
\end{tcolorbox}
\section[moins le quart]{moins le quart — quarter to}

\label{lesson:FR-C08-moins}



\begin{quote}
Twice now you have \textbf{added} to the hour. The last piece of the clock does the opposite, and it does it with a word you can read on a calculator.
\end{quote}

\begin{sounds}
\begin{itemize}
\raggedright
  \item \textbf{moins} = \emph{mwa(n)}. The \textbf{oi} is \emph{wa}, as in \textbf{trois}, and the \textbf{in} is nasal, as in \textbf{cinq}.
  \item The final \textbf{s} is silent. Two familiar sounds and one silent letter — nothing in this word is new.
\end{itemize}
\end{sounds}

\begin{cousinweb}
\textbf{moins} is Latin \textbf{minus}, and it means exactly what \emph{minus} means: \textbf{less}, taken away. English uses the Latin form untouched for subtraction, and keeps the same root in \textbf{minor}, \textbf{minimum} and \textbf{diminish}.
\end{cousinweb}

\begin{grammarlens}[title={Grammar Lens: French subtracts from the next hour}]
To say \textquotedblleft{}quarter to,\textquotedblright{} French does not look back at the hour that has passed. It names the hour \textbf{coming}, and takes the quarter off it:

\begin{quote}
\textbf{Il est trois heures moins le quart.} — \textquotedblleft{}It is a quarter to three.\textquotedblright{}
\end{quote}

Literally: \emph{\textquotedblleft{}it is three hours \textbf{minus} the quarter.\textquotedblright{}} English does the same arithmetic and hides it behind the word \emph{to}; French shows its working.

So the three additions to the hour are one small system:

\begin{itemize}
\raggedright
  \item \textbf{et quart} — the hour, plus a quarter
  \item \textbf{et demie} — the hour, plus a half
  \item \textbf{moins le quart} — the \emph{next} hour, minus a quarter
\end{itemize}
\end{grammarlens}

\subsection*{Your turn}
Say these aloud:

\begin{itemize}
\raggedright
  \item \textquotedblleft{}moins\textquotedblright{} — \emph{mwa(n)}, no \emph{s}
  \item \textquotedblleft{}Il est trois heures moins le quart.\textquotedblright{}
  \item the three in a row — \textquotedblleft{}deux heures et quart, deux heures et demie, trois heures moins le quart\textquotedblright{}
\end{itemize}

\begin{tcolorbox}[breakable,colback=teal!4,colframe=teal!35!black,title={Before you move on}]
How do you say quarter to three? (\emph{Il est trois heures \textbf{moins le quart}.}) Which hour does French name in that sentence, the one gone or the one coming? (\textbf{The one coming} — it subtracts from it.) What Latin word is \emph{moins}? (\textbf{Minus}.) Next: the whole clock, put to use.
\end{tcolorbox}
\section[Practice]{Practice — reading the clock}

\label{lesson:FR-C08-practice}



\begin{quote}
One frame, one noun and three small additions. Put them together and you can say any time on a clock face.
\end{quote}

\subsection*{Your turn: whole hours}
Put each of these into the frame — \emph{Il est …} — then cover the right column and read back from the left:

\noindent
\begin{tabularx}{\linewidth}{@{}>{\raggedright\arraybackslash}X>{\raggedright\arraybackslash}X@{}}
\toprule
\textbf{French} & \textbf{English} \\
\midrule
\textbf{une heure} & one o'clock \\
\textbf{deux heures} & two o'clock \\
\textbf{six heures} & six o'clock \\
\textbf{dix heures} & ten o'clock \\
\bottomrule
\end{tabularx}

\emph{Twice through:} the four whole hours, each inside \emph{Il est …}

Listen for the liaison each time: \emph{ün-eur}, \emph{deu-z-eur}, \emph{see-z-eur}, \emph{dee-z-eur}.

\subsection*{Your turn: the two that take a name}
\begin{itemize}
\raggedright
  \item \textbf{Il est midi.} — noon, \textquotedblleft{}the day's middle\textquotedblright{}
  \item \textbf{Il est minuit.} — midnight, \textquotedblleft{}the night's middle\textquotedblright{}
\end{itemize}

No number and no \emph{heures}: each of them \textbf{is} the hour.

\subsection*{Your turn: the parts of the hour}
\noindent
\begin{tabularx}{\linewidth}{@{}>{\raggedright\arraybackslash}X>{\raggedright\arraybackslash}X@{}}
\toprule
\textbf{French} & \textbf{English} \\
\midrule
\textbf{deux heures et quart} & quarter past two \\
\textbf{deux heures et demie} & half past two \\
\textbf{trois heures moins le quart} & quarter to three \\
\textbf{midi et demi} & half past twelve \\
\bottomrule
\end{tabularx}

Say all four inside the frame and hear the clock move. The third names the hour \textbf{ahead} and subtracts from it.

\begin{grammarlens}[title={What you've built this chapter}]
\begin{itemize}
\raggedright
  \item \textbf{heure} — the same word as English \emph{hour}, both Latin \emph{hōra}, both keeping a silent \emph{h} neither language ever pronounced.
  \item \textbf{il est … heures} — French counts the hours; English says \emph{o'clock}.
  \item \textbf{agreement you cannot hear} — \emph{une heure} against \emph{deux heures}, \emph{et demie} against \emph{et demi}. The page shows it; the ear uses the number in front.
  \item \textbf{midi} and \textbf{minuit} — one idea twice: the \emph{middle} of the day and of the night, with the \emph{-di} of \emph{lundi} and the \emph{nuit} of your first chapter.
  \item \textbf{adding and subtracting} — \emph{et quart}, \emph{et demie}, \emph{moins le quart}.
\end{itemize}
\end{grammarlens}

\begin{tcolorbox}[breakable,colback=teal!4,colframe=teal!35!black,title={Before you move on}]
Say one o'clock, then two. (\emph{Il est une heure. Il est deux heures.}) Say noon and midnight, and what each means literally. (\emph{Midi}, the day's middle; \emph{minuit}, the night's middle.) Say quarter past two, half past two, quarter to three. Which of those names an hour that has not arrived? (\textbf{Quarter to} — \emph{moins le quart} subtracts.) Why do \emph{heure} and \emph{hour} both carry a silent \emph{h}? (Both are Latin \emph{hōra}, where it was already silent.)
\end{tcolorbox}
